\begin{claim}{D.54}{Proved}{No fixed character for a cubic torus action}
Let $a$ have order three in the full
once-punctured-torus loop-outer group.
For every prime $q\ne3$, it fixes no nonzero
continuous character
\[
                         \chi:F\longrightarrow C_q,
                      \qquad F=\widehat F(x,y).         \tag{D.54a}
\]
Consequently the reduction of $A_q(a)$ modulo
$q$ has no eigenvalue one.
For every prime $p$, the endomorphism
$A_p(a)-I$ of $\Z_p^2$ is injective.
The full fixed group of every actual
order-three representative is therefore
contained in the full commutator subgroup
$F'$.

In particular, $A_2$ is injective on every
finite subgroup of the full torus group.
The allowed list of D.52 loses
$C_{12}$ and $D_{24}$.
At a nontrivial cubic multiplier
coordinate $q\ne3$, the remaining
homology matrix is exactly
\[
                         A_q(a)=\lambda(a)_q^2 I.
                                                               \tag{D.54b}
\]
This statement does not exclude the
remaining scalar cubic multiplier case.
\end{claim}
\begin{proof}
Use D.27 to choose an actual representative
$\alpha$ of order three and put
$E=F\rtimes_\alpha\langle s\rangle$,
where $s$ has quotient label $a$.
Write $c=[x,y]$ and $L=\langle c\rangle$.
D.48 supplies a labelled cusp-normalizer
element $v$ with
\[
       v\in N_E(L),\qquad v^3=c^\kappa,\qquad
       vcv^{-1}=c^\lambda,\qquad
                         \kappa_3\in\Z_3^\times.        \tag{D.54c}
\]
It also gives $\lambda_3=1$ and
\[
                  I+A_3(a)+A_3(a)^2=0
                         \quad\text{on }\Z_3^2.         \tag{D.54d}
\]

Suppose $\chi$ in (D.54a) were nonzero
and invariant. Its extension
\[
       \widetilde\chi:E\longrightarrow C_q,\qquad
              \widetilde\chi(hs^j)=\chi(h)
                                                               \tag{D.54e}
\]
is an actual continuous homomorphism.
Put
\[
               N=\ker\chi,\qquad
               E_N=\ker\widetilde\chi=N\rtimes\langle s\rangle .
                                                               \tag{D.54f}
\]
Both have their indicated index $q$
in $F$ and $E$. Since $\chi(c)=0$,
(D.54c) gives
$3\widetilde\chi(v)=0$. The prime $q$
differs from three, so $v\in E_N$.
Thus $v=ns$ for a specified $n\in N$.

The restriction of $\chi$ to the dense
discrete free group is onto $C_q$.
Its covering surface is the degree-$q$
cover of the punctured torus with
boundary monodromy zero.
It has $q$ boundary lifts of degree one,
Euler characteristic $-q$, and hence
type $(1,q)$. Its full fundamental-group
completion is exactly $N$, by the
finite-index completion map.
Choose a discrete word $f$ with
$\chi(f)=1$. Its boundary representatives
are
\[
                c_i=f^i c f^{-i}\quad(0\leq i<q).
                                                               \tag{D.54g}
\]
Their homology classes have exactly the
single relation $\sum_i[c_i]=0$.
In particular each $[c_i]$ is nonzero
in $H_1(N,\F_3)$, because $q\geq2$.
Equivalently, omitting any other boundary
from the usual genus-one presentation
puts $c_i$ into a free basis.

The action on every boundary class has
multiplier $\lambda$ and fixes its label.
This follows from the actual $v\in E_N$:
the element
$n_i=(vf^iv^{-1})f^{-i}$ belongs to $N$,
and
$vc_iv^{-1}=n_i c_i^\lambda n_i^{-1}$.
Conjugation by $s=n^{-1}v$ differs from
this by an inner automorphism of $N$.
Thus $\alpha$ fixes all $q$ boundary
homology classes over $\Z_3$.

We give the required splitting of
the entire homology module, including
its comparison with the base torus.
Let $D=F/N=C_q$ be the deck group and set
\[
          M=H_1(N,\Z_3),\qquad
                       e=\frac1q\sum_{d\in D}d .
                                                               \tag{D.54h}
\]
The action of $\alpha$ on $M$ commutes
with the deck action: it induces the
identity on $D$, and two lifts of a
deck transformation differ by an inner
automorphism of $N$ on this homology.
Consequently $e$ commutes with $\alpha$.

Filling the $q$ boundaries gives the
compact degree-$q$ torus cover.
Explicitly its lattice is the kernel
$\Lambda\leq\Z^2$ of the character
$\Z^2\to C_q$ induced by $\chi$,
and its map to the base torus is
$\R^2/\Lambda\to\R^2/\Z^2$.
The deck maps are translations and
act trivially on the filled homology.
The boundary submodule
\[
              B_\partial
                =\Z_3[D]/\Z_3\Bigl(\sum_{d\in D}d\Bigr)
                                                               \tag{D.54i}
\]
has no deck-invariant vector, because
$q$ is a unit in $\Z_3$.
The exact boundary sequence and (D.54h)
therefore give
\[
             (1-e)M=B_\partial,\qquad
             eM\cong H_1(\R^2/\Lambda,\Z_3)
                         \cong H_1(\R^2/\Z^2,\Z_3).
                                                               \tag{D.54j}
\]
The last isomorphism is induced by the
displayed covering: its lattice inclusion
has index $q$, a $3$-adic unit.
These maps intertwine the induced
automorphisms, so $\alpha|_{eM}$ is
exactly $A_3(a)$.
On $(1-e)M$ it is the identity by
(D.54g). Hence
\[
                  I+\alpha+\alpha^2=0
                       \quad\text{on }H_1(N,\F_3):
                                                               \tag{D.54k}
\]
on the first summand this is (D.54d),
and on the boundary summand the norm
is multiplication by three.

The cube of the actual element $ns$
now gives the contradiction:
\[
 \begin{split}
  [v^3]
      &=[n\,\alpha(n)\,\alpha^2(n)]\\
      &=(I+\alpha+\alpha^2)[n]=0
                       \quad\text{in }H_1(N,\F_3),\\
  [c^\kappa]&=(\kappa_3\bmod3)[c]\ne0 .
 \end{split}                                             \tag{D.54l}
\]
Equations (D.54c) and (D.54l) concern
the same two elements of $N$.
Thus no invariant character (D.54a)
exists.

For $q\ne3$, a matrix and its transpose
have the same rank, so a one-eigenvector
would supply a nonzero invariant
character. It has just been excluded.
In particular $A_q(a)-I$ is invertible
over $\Z_q$.
At $q=3$ its determinant is three,
by D.48, so it is still injective.
If $h$ is fixed by an actual representative,
its image in every $\Z_p^2$ is therefore
zero. The common kernel of the actual
abelianization map
$F\to\widehat\Z^2=\prod_p\Z_p^2$
is $F'$, proving the fixed-group assertion.

At a nontrivial cubic multiplier
coordinate the two eigenvalue multisets
were exactly
$\{1,\lambda_q\}$ and
$\{\lambda_q^2,\lambda_q^2\}$ in D.51.
The first is now excluded.
A finite-order matrix in characteristic
zero is diagonalizable, so the second
gives the actual scalar matrix
(D.54b).

Finally let $G$ be any finite torus
subgroup. D.33 gives $|G|\mid24$.
A nontrivial kernel of $A_2|_G$ would
contain a prime-order element.
Its prime can only be two or three.
D.33 excludes an element of order two,
and (D.54a), with $q=2$, excludes
an element of order three acting
trivially on two-adic homology.
Thus $A_2|_G$ is injective.
The faithful-image part (D.52h), together
with D.50, now excludes $C_{12}$ and
$D_{24}$ from that allowed list.

The program
\path{code/part_c/torus_cover_norm_d54_d55.py}
builds the cyclic covering homology,
its individual boundary vectors and
the deck-average splitting for several
primes $q\ne3$, and verifies the forced
norm contradiction over $\F_3$.
The displayed maps prove it for every
such prime.
\end{proof}
